\section{Experimental Setup}

We design experiments to answer two sequential questions:

\textbf{RQ1:} Do spurious features dominate core features early in ERM training? (Existence)

\textbf{RQ2:} Do spurious features receive stronger gradient signals than core features? (Mechanism)

\subsection{Dataset}

\textbf{Waterbirds v1.0} \citep{sagawa2020distributionally}: A benchmark for studying spurious correlations in image classification. The task is to classify bird type (waterbird vs.\ landbird) while the background habitat (water vs.\ land) serves as a spurious feature correlated with the label in training data.

\begin{table}[h]
\centering
\begin{tabular}{@{}ll@{}}
\toprule
Statistic & Value \\
\midrule
Total images & 11,788 \\
Training & 4,795 \\
Validation & 1,199 \\
Test & 5,794 \\
Classes & 2 (waterbird, landbird) \\
Spurious attribute & Background (water, land) \\
Minority groups & Waterbird-land, Landbird-water \\
\bottomrule
\end{tabular}
\caption{Waterbirds dataset statistics.}
\label{tab:dataset}
\end{table}

\subsection{Model and Training}

\begin{table}[h]
\centering
\begin{tabular}{@{}ll@{}}
\toprule
Parameter & Value \\
\midrule
Architecture & ResNet-50 (ImageNet pretrained) \\
Optimizer & SGD (momentum=0.9, weight\_decay=1e-4) \\
Learning rate & 0.01 (H-E1), 0.001 (H-M1) \\
LR schedule & StepLR (step=20, gamma=0.1) \\
Batch size & 128 \\
Epochs & 50 \\
Seeds & 1 (H-E1), 3 (H-M1) \\
\bottomrule
\end{tabular}
\caption{Training configuration.}
\label{tab:training}
\end{table}

\subsection{Implementation}

Experiments implemented in PyTorch 2.0 with pytorch-grad-cam for attribution. Training on single NVIDIA GPU. Gradient norms computed via \texttt{torch.autograd.grad} with per-sample separation. All code available in supplementary material.
